Seja a soma dos quadrados dos residuos da reta ajustada
\[
S(m,b)=\sum_{i=1}^{n}(y_i-mx_i-b)^2.
\]
As formas finais das eqs. (4.51) e (4.52) surgem impondo as condicoes de otimo
\[
\frac{\partial S}{\partial m}=0,
\qquad
\frac{\partial S}{\partial b}=0.
\]

Derivando em relacao a $m$,
\[
\frac{\partial S}{\partial m}=-2\sum_{i=1}^{n}x_i(y_i-mx_i-b)=0,
\]
o que fornece
\[
\sum_{i=1}^{n}x_i y_i=m\sum_{i=1}^{n}x_i^2+b\sum_{i=1}^{n}x_i.
\]
Essa e a forma final de uma das equacoes normais.

Derivando em relacao a $b$,
\[
\frac{\partial S}{\partial b}=-2\sum_{i=1}^{n}(y_i-mx_i-b)=0,
\]
obtemos
\[
\sum_{i=1}^{n}y_i=m\sum_{i=1}^{n}x_i+nb.
\]

Portanto, as eqs. (4.51) e (4.52) sao exatamente as duas equacoes normais do ajuste linear por minimos quadrados, obtidas sem passos ocultos.